% authority: science_draft
\documentclass[11pt]{article}

\usepackage[margin=1in]{geometry}
\usepackage{amsmath,amssymb,amsthm}
\usepackage{booktabs}
\usepackage{enumitem}
\usepackage{longtable}

\newtheorem{definition}{Definition}
\newtheorem{theorem}{Theorem}
\newtheorem{proposition}{Proposition}
\newtheorem{remark}{Remark}
\newtheorem{nonconclusion}{Non-conclusion}

\newcommand{\Ksrc}{K_{\mathrm{src}}}
\newcommand{\Aint}{A_{\mathrm{int}}}
\newcommand{\Bdef}{B_{\mathrm{def}}}
\newcommand{\Gint}{G_{\mathrm{int}}}
\newcommand{\Rdecl}{\widehat G_{\mathrm{decl}}}
\newcommand{\Conf}{\mathcal C_{\mathrm{mat}}}
\newcommand{\Sect}{\mathrm{Sec}_{\mathrm{src}}}
\newcommand{\Malformed}{\bot_{\mathrm{mat}}}
\newcommand{\Undetermined}{\mathsf{undetermined}_{\mathrm{mat}}}
\newcommand{\MetricData}{\mathrm{MetricData}}
\newcommand{\geff}{g_{\mathrm{eff}}}

\title{P7-T01 Source-Matter Ontology and Sector Taxonomy v1}
\author{Ontology Formalizer Packet RT-20260727-007}
\date{2026-07-27}

\begin{document}
\maketitle

\section{Control status and result}

This artifact is the fused old-style science-draft output of
\texttt{AJ-RT-20260727-007-001}. It constructs
\[
  \mathrm{SourceMatterIncidenceCandidate}_{v1}
\]
as a \emph{draft/control, proposal-only} source-side ontology candidate. The
candidate supplies explicit configuration variables, internal labels, source
charge and current chains, defect classes, sector maps, source-invariant
bookkeeping records, a nontrivial finite configuration, a null
configuration, a malformed branch, and one exact naturality theorem within
the candidate.

The artifact does not edit canonical ontology or adopt a source law. It does
not establish physical matter, physical charge, a physical conservation law,
detector semantics, an effective geometry, a target metric, stress energy, a
matter action, matter coupling, Einstein equations, benchmark promotion,
proof authority, publication, completed derivation, a global no-go theorem,
or future source-extension impossibility. Current adoption remains
\texttt{blocked\_adoption\_open\_continuation}.

\section{Why a new candidate is required}

The registered legacy matter records provide useful target grammar,
matter-sector labels, source probe and readout syntax, compatibility
certificates, candidate-law obligations, and scoped evidence/preconditions.
Their own status controls expressly prevent those interfaces from being read
as a source-derived matter ontology. In particular, a label called a matter
sector is not yet an intrinsic degree of freedom, and a passing certificate is
not a law that produces a sector.

P7-T01 therefore asks a narrower mathematical question:

\begin{quote}
Can one write a concrete source-only configuration space and sector map that
does not assume the frozen P6 effective-geometry output, a target metric,
stress energy, a matter action, coupling, detector semantics, or GR behavior?
\end{quote}

The construction below answers that question at proposal-only candidate
level. It does not answer whether the project should adopt this candidate.

\section{Declared source data}

\begin{definition}[Finite augmented source incidence complex]
A finite augmented source incidence complex is a chain complex of free
abelian groups
\[
 C_2(\Ksrc;\mathbb Z)
 \xrightarrow{\partial_2}
 C_1(\Ksrc;\mathbb Z)
 \xrightarrow{\partial_1}
 C_0(\Ksrc;\mathbb Z)
 \xrightarrow{\epsilon}
 \mathbb Z,
 \qquad
 \partial_1\partial_2=0,
 \qquad
 \epsilon\partial_1=0,
\]
with a declared finite vertex basis \(V\), oriented source-link basis \(E\),
and source-relation basis \(F\). The boundary matrices are source data.
Each source-link boundary has one declared positive and one declared negative
endpoint; a self-loop may have zero boundary. The augmentation sends every
vertex basis element to \(1\).
``Incidence'' here is combinatorial source syntax. It is not a target atlas,
spacetime topology, metric, causal cone, proper-time relation, or physical
locality claim.
\end{definition}

\begin{definition}[Internal candidate data]
The candidate declares:
\begin{enumerate}[leftmargin=*]
\item a finitely generated abelian group
  \(\Aint\), used only as a source charge-label group;
\item a finitely generated abelian group
  \(\Bdef\), used only as a source defect-label group;
\item an internal group \(\Gint\) and a finite declared set
  \(\Rdecl\) of representation isomorphism classes, each with a declared
  carrier and homomorphism
  \(\varrho_\rho:\Gint\to\mathrm{GL}(W_\rho)\);
\item a distinguished empty representation label \(0_R\);
\item a source charge-label map
  \[
    \chi:\Rdecl\cup\{0_R\}\longrightarrow\Aint,
    \qquad \chi(0_R)=0.
  \]
\end{enumerate}
These are proposal-only internal labels. They are not measured charges,
particle species, Standard Model representations, detector outcomes, or
empirical calibrations.
\end{definition}

\begin{definition}[Model-to-world status ladder]
The packet keeps four levels separate:
\begin{enumerate}[leftmargin=*]
\item the incidence complex, coefficient groups, internal group,
  representations, \(\chi\), occupations, and transition grammar are
  \emph{stipulated candidate primitives};
\item charge chains, representation grades, sector coordinates, current
  constraints, and defect classes are \emph{candidate-derived mathematical
  records};
\item reference to physical matter, charge, motion, measurement, or
  spacetime is \emph{open and not supplied};
\item canonical ontology or source-law adoption is \emph{protected and
  human-gated}.
\end{enumerate}
An exact consequence at the second level does not move an object to the
third or fourth level.
\end{definition}

Coefficient extension gives chain groups
\(C_i(\Ksrc;\Aint)\) and \(C_i(\Ksrc;\Bdef)\). Let
\[
  \epsilon_{\Aint}:C_0(\Ksrc;\Aint)\longrightarrow\Aint,
  \qquad
  \epsilon_{\Aint}\!\left(\sum_{v\in V}a_v[v]\right)
  =\sum_{v\in V}a_v .
\]
For an incidence boundary, every oriented link contributes one \(+\) and one
\(-\) endpoint, so
\(\epsilon_{\Aint}\partial_1=0\).

\section{Source-matter configurations}

\begin{definition}[Proposal-only source-matter configuration]
A source-matter configuration has primitive tuple
\[
  \Psi=(m,z)
\]
with
\[
  m:V\times\Rdecl\longrightarrow\mathbb N_0,
  \qquad
  z\in Z_1(\Ksrc;\Bdef)=\ker\partial_1 .
\]
Here \(m\) is the source representation-occupation multiplicity and \(z\) is
a closed source defect chain. Occupation, local charge, charge chain, and
representation grade are derived:
\[
  \nu(v)=\sum_{\rho\in\Rdecl}m(v,\rho),\qquad
  q_v=\sum_{\rho\in\Rdecl}m(v,\rho)\chi(\rho),
\]
\[
  q(\Psi)=\sum_{v\in V}q_v[v],
  \qquad
  M_\rho(\Psi)=\sum_{v\in V}m(v,\rho).
\]
It is admissible when:
\begin{enumerate}[label=\textup{(C\arabic*)},leftmargin=*]
\item all vertices, links, relations, coefficient groups, representation
  carriers, actions, labels, and values of \(\chi\) are declared;
\item \(m(v,\rho)\) is nonnegative on the declared finite domain, and
  \(\nu,q,M\) are the derived fields above rather than independent tags;
\item \(z\) is a declared one-chain and \(\partial_1z=0\);
\item no component is defined from a target atlas, target metric, effective
  geometry, stress-energy variable, matter action, detector outcome,
  benchmark behavior, validator result, registry row, role authority,
  generated derivative, local cache, or commit state.
\end{enumerate}
Failure of any condition returns \(\Malformed\) with the failed condition and
offending datum recorded.
\end{definition}

\begin{definition}[Candidate state space]
The proposal-only state space is
\[
  \Conf(\Ksrc;\Aint,\Bdef,\Gint,\Rdecl)
  =
  \{\Psi:\Psi\text{ satisfies (C1)--(C4)}\}
  \cup\{\Malformed\}.
\]
This is a kinematic candidate. It supplies no action, Hamiltonian, evolution
operator, probability rule, energetic weighting, or physical time.
\end{definition}

\begin{definition}[Source sectors, representation grade, and invariant records]
For admissible \(\Psi=(m,z)\), define
\[
  Q(\Psi)=\epsilon_{\Aint}(q(\Psi)),\qquad
  D(\Psi)=[z]\in
  H_1(\Ksrc;\Bdef)=
  \ker\partial_1/\operatorname{im}\partial_2,
\]
and use the derived \(M(\Psi)\) as a representation grade. The
transition-invariant candidate sector map is
\[
  \Sect(\Psi)=\bigl(Q(\Psi),D(\Psi)\bigr).
\]
Also define the source-invariant bookkeeping record
\[
  O_{\mathrm{src}}(\Psi)
  =
  \left(
    N(\Psi)=\sum_v\nu(v),\,
    Q(\Psi),\,
    D(\Psi),\,
    M(\Psi)
  \right).
\]
The grade \(M\) is equivariant under internal-label transport but is not
claimed to be transition-invariant unless a later packet constrains the
update record \(\Pi\) accordingly.
These are source invariants or labels. They are not empirical or operational
observables. Invariance licenses no observer, detector, measurement,
calibration, unit, localization, or operational-geometry semantics.
\end{definition}

\begin{definition}[Identity ladder and underdetermined result]
Four levels must remain distinct:
\begin{enumerate}[leftmargin=*]
\item \emph{presentation identity}: equality of every declared field;
\item \emph{candidate isomorphism}: transport by the declared source-chain
  and internal-label isomorphisms;
\item \emph{sector coincidence}: equality, or declared transport, of
  \(\Sect\), which is coarser than configuration identity;
\item \emph{physical identity}: not supplied by this packet.
\end{enumerate}
Candidate isomorphism is not the unresolved general
\(\mathrm{EqSrc}\) relation, observational equivalence, or physical gauge.
If all fields are well typed but the candidate rules do not decide a requested
classification, the evaluator returns \(\Undetermined\) with the missing rule
recorded. It does not replace missing information by zero, the valid null
state, or \(\Malformed\).
\end{definition}

\section{Currents and kinematically admissible transitions}

\begin{definition}[Source-current transition record]
For admissible endpoint configurations
\(\Psi^-=(m^-,z^-)\) and
\(\Psi^+=(m^+,z^+)\), with their derived charge chains \(q^-\) and
\(q^+\), a candidate transition record is
\[
  T=(\Psi^-,J,w,\Pi,\Psi^+)
\]
where \(J\in C_1(\Ksrc;\Aint)\) is a source current-label chain,
\(w\in C_2(\Ksrc;\Bdef)\) is a defect-relation chain, and \(\Pi\) is the
declared source-only update from \(m^-\) to \(m^+\).
It is kinematically admissible when
\begin{align}
  q^+-q^- &= \partial_1J,
    \label{eq:continuity}\\
  z^+-z^- &= \partial_2w,
    \label{eq:defect}
\end{align}
and \(\Pi\) uses only declared source data. Equation
\eqref{eq:continuity} is a candidate chain constraint, not an adopted
dynamical or physical conservation law. P7-T01 does not select which
admissible transition occurs.
\end{definition}

\begin{proposition}[Candidate total-charge and defect-sector preservation]
For every kinematically admissible transition,
\[
  Q(\Psi^+)=Q(\Psi^-),\qquad
  D(\Psi^+)=D(\Psi^-).
\]
\end{proposition}

\begin{proof}
Apply \(\epsilon_{\Aint}\) to
\(q^+-q^-=\partial_1J\). Since
\(\epsilon_{\Aint}\partial_1=0\), the total source charge labels agree.
Equation \eqref{eq:defect} says \(z^+\) and \(z^-\) differ by an element of
\(\operatorname{im}\partial_2\), so they define the same homology class.
\end{proof}

\begin{remark}[Conditional consistency rather than independent dynamics]
The preceding proposition is a conditional chain-level invariant. Because
equations \eqref{eq:continuity} and \eqref{eq:defect} are part of the
definition of a kinematically admissible transition, the proposition checks
consistency of that proposal. It is not an independently derived equation of
motion, Noether theorem, or physical conservation law.
\end{remark}

\section{Sector and charge taxonomy}

\begin{longtable}{p{0.19\linewidth}p{0.28\linewidth}p{0.43\linewidth}}
\toprule
Candidate class & Defining datum & Exact licensed meaning \\
\midrule
\endhead
Null configuration stratum &
\(m=0,\ z=0\), hence \(N=q=M=0\) &
Empty source configuration. Other nontrivial configurations can share its
\((Q,D)=(0,[0])\) sector, so nullity is not a sector property. It is not a
physical vacuum or ground state.\\
Charge sector &
\(Q\in\Aint\) &
Additive source label obtained by augmentation. It is not electric, color,
baryon, energy, or another physical charge without later authority.\\
Defect sector &
\([z]\in H_1(\Ksrc;\Bdef)\) &
Equivalence class of a closed combinatorial source chain modulo declared
source relations. It is not a spacetime topological defect.\\
Representation grade &
\(M:\Rdecl\to\mathbb N_0\) &
Multiplicity of declared internal representation labels. It is not a
particle-species or empirical classification and need not be transition
invariant.\\
Invariant sector &
\((Q,[z])\) &
Product candidate sector preserved by the declared kinematic constraints.\\
Current record &
\(J\in C_1(\Ksrc;\Aint)\) with
\(q^+-q^-=\partial_1J\) &
Chain-level transport witness between declared endpoint labels. It is not
physical flux or stress energy.\\
Malformed branch &
\(\Malformed\) plus failed condition &
Fail-closed result for an undeclared object, support violation, noncycle,
continuity failure, defect-relation failure, or forbidden import.\\
\bottomrule
\end{longtable}

\section{Internal symmetry and source relabeling}

\begin{definition}[Admissible change of source presentation]
An admissible change of source presentation from candidate data
\((\Ksrc,\Aint,\Bdef,\Gint,\Rdecl)\) to primed data consists of:
\begin{enumerate}[leftmargin=*]
\item chain-group isomorphisms
  \(F_i:C_i(\Ksrc;\mathbb Z)\to C_i(\Ksrc';\mathbb Z)\) satisfying
  \(F_{i-1}\partial_i=\partial_i'F_i\), with \(F_0\) induced by a vertex
  bijection;
\item group isomorphisms
  \(\alpha:\Aint\to\Aint'\) and
  \(\beta:\Bdef\to\Bdef'\);
\item an internal-group isomorphism
  \(\phi:\Gint\to\Gint'\) and compatible bijection
  \(\gamma:\Rdecl\to\Rdecl'\), transporting representation classes so that
  \(\varrho'_{\gamma\rho}\circ\phi\) is isomorphic to
  \(\varrho_\rho\), and satisfying
  \[
    \chi'(\gamma\rho)=\alpha(\chi(\rho));
  \]
\item coefficientwise chain maps
  \(L_i^A=F_i\otimes\alpha\) and
  \(L_i^B=F_i\otimes\beta\), together with occupation transport
  \(m'(F_0v,\gamma\rho)=m(v,\rho)\) and fieldwise transport of
  \(z,J,w,\Pi\).
\end{enumerate}
The admitted transformations are exactly this compatible subgroup of the
product of source-chain, coefficient-group, internal-group, and label
isomorphisms; arbitrary independent permutations are not admitted.
This is source-presentation naturality only. It is not target-coordinate
covariance, diffeomorphism invariance, local gauge symmetry, or physical
equivalence.
\end{definition}

An internal automorphism need not fix the raw label \(Q\). The theorem below
therefore claims equivariant transport \(Q\mapsto\alpha Q\), not invariance of
the raw value. A later use that calls such a transformation a
within-sector symmetry must either restrict to \(\alpha Q=Q\) or replace the
charge coordinate by its explicitly declared automorphism orbit.

For example, take \(\Aint=\mathbb Z\), labels
\(\rho_+,\rho_-\) with \(\chi(\rho_\pm)=\pm1\), the automorphism
\(\alpha(n)=-n\), and let \(\gamma\) exchange the two labels. Then
\(\chi'\gamma=\alpha\chi\) and \(Q\mapsto-Q\). This is exact equivariance,
not literal invariance of \(Q\).

\begin{theorem}[P7T01-THM-SOURCE-RELABELING-SECTOR-NATURALITY-001]
Let \((F,\phi,\alpha,\beta,\gamma)\) be an admissible change of source
presentation. Then:
\begin{enumerate}[leftmargin=*]
\item admissible configurations transport to admissible configurations;
\item kinematically admissible transitions transport to kinematically
  admissible transitions;
\item the sector map and representation grade are equivariant:
  \[
    \Sect(\Psi')=
    \left(
      \alpha Q(\Psi),\,
      H_1(F;\beta)D(\Psi)
    \right),
    \qquad
    M(\Psi')=\gamma_*M(\Psi);
  \]
\item a continuity violation remains a continuity violation under the
  isomorphism, so that malformed branch cannot be repaired by relabeling.
\end{enumerate}
\end{theorem}

\begin{proof}
The vertex and representation-label bijections preserve occupation
multiplicity. Compatibility of \(\chi\) gives, at every vertex,
\[
  q'_{F_0v}
  =
  \sum_\rho m(v,\rho)\chi'(\gamma\rho)
  =
  \alpha\!\left(\sum_\rho m(v,\rho)\chi(\rho)\right)
  =
  \alpha(q_v).
\]
Thus charge is transported as a derived field, not as an independently
relabelled tag. The coefficientwise chain-map identity gives
\[
  \partial_1'L_1^A J
  =
  L_0^A\partial_1J
  =
  L_0^A(q^+-q^-),
\]
hence the continuity condition is preserved. The same argument using
\(L_1^B\partial_2=\partial_2'L_2^B\) preserves the defect-relation
condition and induces \(H_1(F;\beta)\).

Vertex bijectivity implies
\(\epsilon_{\Aint'}L_0^A=\alpha\epsilon_{\Aint}\), yielding the stated
charge component. Representation grades are permuted by \(\gamma\).
Finally, if
\(q^+-q^--\partial_1J\neq0\), its image under the coefficientwise
isomorphism \(L_0^A\) is nonzero, so the corresponding primed transition
still fails. No target structure enters the argument.
\end{proof}

\begin{remark}[Status of the theorem]
The proof is exact for the declared finite candidate. Its conclusion is
mathematical naturality of proposal-only source records. It is not a theorem
that nature realizes the candidate, not a physical conservation law, not an
ontology adoption, and not proof authority for downstream GR claims.
\end{remark}

\section{Explicit configurations}

\subsection{Nontrivial finite witness}

Let \(V=\{u,v\}\). Let \(C_1\) have source-link basis
\(\{e,\ell\}\), with
\[
  \partial_1 e=[v]-[u],
  \qquad
  \partial_1\ell=0,
\]
and let \(C_2=0\). Take
\(\Aint=\Bdef=\mathbb Z\), declare an internal representation label
\(\rho\in\Rdecl\) with \(\chi(\rho)=1\), and define:
\[
\begin{array}{c|cccc}
 & m & \nu & q & z\\
\hline
\Psi^- & m(u,\rho)=1 & \nu(u)=1 & [u] & \ell\\
\Psi^+ & m(v,\rho)=1 & \nu(v)=1 & [v] & \ell
\end{array}
\]
with all unlisted multiplicities zero. The displayed \(\nu\) and \(q\) are
derived from \(m\) and \(\chi\). Put
\[
  J=e,\qquad w=0,
\]
and let \(\Pi\) record the source-label move \(u\to v\). Then
\[
  q^+-q^-=[v]-[u]=\partial_1e=\partial_1J,
  \qquad
  z^+-z^-=0=\partial_2w.
\]
Both endpoints therefore occupy the nontrivial invariant sector
\[
  \bigl(1,[\ell]\bigr)
\]
and have representation grade \(M_\rho=1\).
Because \(C_2=0\) and \(\ell\neq0\), \([\ell]\) is nonzero. This is a
finite source-chain witness, not a particle trajectory, spacetime loop,
physical worldline, or measured charge flow.

\subsection{Null configuration}

The tuple
\[
  \Psi_0=(m=0,z=0)
\]
is admissible, has derived \(\nu=q=M=0\), and has sector \((0,[0])\).
It is a null source record, not a physical vacuum, minimum-energy state, or
cosmological background.

\subsection{Nontrivial neutral configuration}

Declare labels \(\rho_+,\rho_-\) with
\(\chi(\rho_+)=1\) and \(\chi(\rho_-)=-1\). On two occupied vertices, take
\[
  m(u,\rho_+)=m(v,\rho_-)=1,\qquad z=0.
\]
Then \(q=[u]-[v]\), \(Q=0\), \(N=2\), and both representation grades are
nonzero. This valid configuration is
not the null configuration. Zero total source charge never implies nullity,
physical neutrality, or absence of matter.

\subsection{Malformed transition}

Keep the two nontrivial endpoints above but set \(J=0\). Then
\[
  q^+-q^-=[v]-[u]\neq0=\partial_1J.
\]
The transition fails equation \eqref{eq:continuity} and returns
\[
  \Malformed(
    \mathrm{failed\_condition}=\mathrm{C\mbox{-}TRANS\mbox{-}CONTINUITY},
    \mathrm{residual}=[v]-[u]).
\]
No target metric, detector output, or validator result may repair it.

\subsection{Underdetermined request}

For a well-typed admissible configuration, a question such as ``which
detector outcome corresponds to this sector?'' has no rule in P7-T01. The
result is \(\Undetermined\) with
\(\mathrm{missing\_rule}=\mathrm{detector\_semantics}\), not the valid null
sector and not \(\Malformed\). The missing rule remains a later burden.

\section{Matter identity and emergence boundary}

Within this candidate, a matter token is individuated only at the
candidate-relative level by a source configuration and its declared
candidate-isomorphism orbit, not by a persistent observer name. Sector
coincidence is weaker than configuration identity and supplies neither
diachronic particle identity nor a superselection theorem.

The packet stipulates source-side matter degrees of freedom. It does not
decide whether physical matter is fundamental, emergent, or a substrate
excitation. Calling a configuration a physical excitation would additionally
require a derived reference state, lawful dynamics, stability, and an
interpretation map, none of which P7-T01 supplies. Canonical physical identity,
empirical detectability, and correspondence to known particles remain open
and human-gated where ontology adoption would be required.

\section{Legacy-evidence mapping without status upgrade}

\begin{longtable}{p{0.25\linewidth}p{0.31\linewidth}p{0.34\linewidth}}
\toprule
Legacy object & Candidate-side mapping & Preserved status and block \\
\midrule
\endhead
\(\mathrm{SourceMatterSector}_{v1}\) &
Possible interface projection of \(\Sect(\Psi)\) plus separately labelled
configuration strata or representation grades. &
Target grammar only; not adopted sector semantics.\\
\(\mathrm{SourceProbeToken}_{v1}\),
\(\mathrm{SourceDetectorReadout}_{v1}\),
\(\mathrm{MatterResponseToken}_{v1}\) &
Excluded from the ontology core; possible later operational-protocol inputs. &
No observer, detector, measurement, or empirical semantics.\\
\(\mathrm{MatSrcDiscLaw}_{v1}\) and
\(\mathrm{SemMatter}_{\mathrm{src}}\) &
The \(\chi\)-linked occupation model and \(\Sect\) are a concrete
proposal-only nonbottom discriminator candidate that can be tested against
the old obligations. &
No source-law or matter-semantics adoption; prior Refuter obstruction is not
retroactively erased.\\
\(\mathrm{MatterSemanticsCandidate}^{\mathrm{cand}}_{v1}\) &
Earlier interface/certificate candidate; may receive a later projection from
the new configuration space. &
No replacement, adoption, detector semantics, or coupling.\\
\(\mathrm{PositiveMSProfile}_{v1}\),
stable-partition and stable-bridge records &
Scoped evidence/preconditions for later compatibility review. &
Remain scoped evidence/preconditions; do not instantiate this candidate.\\
\(\mathrm{SourceMatterSemanticsAdoptionReadinessLaw}_{v1}\) &
Readiness-context cross-check only. &
Scoped evidence/precondition; the word ``Law'' does not adopt a law.\\
\(\mathrm{RR}_E\) and source-certificate records &
Declared-scope transport and fail-closed record grammar only. &
No detector equivalence, matter semantics, or proof authority.\\
\(\mathrm{PTMI}_E\) and
\(\mathrm{MatterCouplingBridgeTarget}_{v1}\) &
Downstream bridge target that may later consume a separately authorized
ontology projection. &
No stress-energy, matter action, coupling law, or matter coupling.\\
P6 Gate B package &
No input to the construction except the prohibition on treating it as
effective geometry. &
Gate B remains zero of eight; the local freeze is preserved.\\
\bottomrule
\end{longtable}

\section{Assumption and dependency summary}

\begin{enumerate}[leftmargin=*]
\item The finite augmented chain complex, coefficient groups, internal group,
  representation carriers, and charge-label map \(\chi\) are explicit
  proposal-only inputs rather than source-derived physical primitives.
\item Boundary compatibility
  \(\partial_1\partial_2=0\) and incidence augmentation
  \(\epsilon\partial_1=0\) are mathematical assumptions of this candidate.
\item No action or evolution rule chooses among admissible transitions.
\item No target atlas, target metric, \(\MetricData(E)\), \(\geff\), physical
  cone, clock, rod, stress tensor, matter action, detector protocol, coupling
  law, or GR behavior is assumed.
\item Source equivalence in this packet is limited to declared chain
  isomorphisms and internal-label isomorphisms; it is not the general
  \(\mathrm{EqSrc}\) relation and physical equivalence is not inferred.
\item Existing matter-profile records retain their exact registered statuses.
\item Any adoption or canonical ontology edit requires separate protected
  human Gate Chair authority.
\end{enumerate}

\section{Distance-to-GR and freeze evaluation}

\[
\begin{array}{ll}
\text{source matter configuration target}
  & \text{proposal-only candidate supplied}\\
\text{sector/charge/current/defect taxonomy}
  & \text{proposal-only finite definitions supplied}\\
\text{canonical matter ontology}
  & \text{not adopted}\\
\text{physical charge or conservation}
  & \text{not established}\\
\text{operational detector semantics}
  & \text{missing}\\
\text{matter action or evolution law}
  & \text{missing; P7-T02 remains separate}\\
\text{stress energy and universal coupling}
  & \text{missing}\\
\text{effective geometry}
  & \text{not assumed; P6 Gate B freeze unchanged}\\
\text{Einstein equations}
  & \text{not derived}\\
\text{exact-GR benchmark}
  & \text{unchanged target-side benchmark}\\
\text{adoption status}
  & \texttt{blocked\_adoption\_open\_continuation}
\end{array}
\]

This candidate is materially distinct from the frozen P6 effective-geometry
family and does not reopen it. Repeating the same incidence candidate without
a new source law, theorem, variation class, stress result, or protected
decision should be locally frozen. A fresh P7-T02 packet may proceed only
after checkpoint and only to construct an explicit source action, evolution
equation, or transition law; this artifact does not preselect one.

\section{Exact-GR recovery obligations}

Any later use of this candidate must:
\begin{enumerate}[leftmargin=*]
\item keep exact GR as a target benchmark rather than a source premise;
\item derive any effective geometry from source structure independently of
  the matter labels defined here;
\item separately derive physical interpretation, operational protocols,
  matter dynamics, stress energy, universal coupling, and Einstein equations;
\item preserve the distinction between task-local mathematical truth,
  canonical ontology status, physical interpretation, and benchmark status;
\item obtain protected human authority before ontology or benchmark
  promotion.
\end{enumerate}

\section{Non-conclusions}

\begin{nonconclusion}
This artifact is not canonical ontology, not a source-law adoption, not a
physical matter ontology, not a particle model, not a physical charge or
conservation law, not an observer or detector theory, not a target atlas,
metric, causal cone, clock, rod, stress-energy tensor, matter action, coupling
law, matter-coupling derivation, Einstein-equation derivation, benchmark
promotion, proof publication, completed derivation, global theory rejection,
or proof that future source extensions are impossible.
\end{nonconclusion}

\section*{Registered source basis}

\begin{itemize}[leftmargin=*]
\item The AEther-Flow Research Project. (2026, July 20).
  \emph{Recommendations implementation plan continue task v21}
  [Internal control plan].
\item The AEther-Flow Research Project. (2026, June 28).
  \emph{Matter-Coupling Bridge Target v1 formalization}
  [Registered draft/control TeX artifact].
\item The AEther-Flow Research Project. (2026, July 2).
  \emph{Source certificate operation laws v1}
  [Registered draft/control TeX artifact].
\item The AEther-Flow Research Project. (2026, June 29--30).
  \emph{Matter-semantics candidate-law, target, candidate, and Refuter
  records} [Registered draft/control TeX artifacts].
\item The AEther-Flow Research Project. (2026, July 27).
  \emph{P6-T08 Gate B review and P15-T04 source-dynamics reconstruction
  report} [Tracked research-control records].
\end{itemize}

\end{document}
